%% title: MTSMs in Python
%% category: theories

\begin{abstract}
A working implementation of the MTSM: a transducer that values each mother node from its two daughters, one feature tier at a time, plus two evaluators that check whether a tree could have been built that way. Every code block on this page can be run in the browser.
\end{abstract}
\label{ch:mtsm}

\section{The implementation}\label{sec:mtsm-code}

The code below builds trees bottom-up from a stream of left daughters. Each mother's value on each tier (person, number, gender) is computed by a transition function $\tau$ from its daughters' values on that tier alone. Fold any indented block with the arrows in the gutter, or press \textbf{run} to load it into a Python session for this page.\footnote{Python runs entirely in your browser, using Pyodide. The first run downloads it, about 10\,MB; after that it is cached.}

\begin{minted}{python}
### MTSM implementation in Python

tier_list = ['person', 'number', 'gender'] 

# A syntactic node. 
class snode:
    def __init__(self, left=None, right=None, mother=None):
        self.vals = {} # list of tier values
        for tier in tier_list:
            self.vals[tier] = None
        self.left = left # left daughter
        self.right = right # right daughter
        self.mother = mother # mother
    def __repr__(self):
        return f"snode(vals={self.vals!r})"

# A transducer which gives values to non-terminals, given a transition function. 
def run_mtsm(input_stream, tau):
    output_stream = []
    right = None; # assumes the bottommost right node is a lexical root unvalued for any features.
    for left in reversed(input_stream):
        mother = tau(left, right) # tau is the transition function
        output_stream.insert(0, right)
        output_stream.insert(0, left)
        right = mother
    output_stream.insert(0, right)
    return output_stream

    # Input: list of left nodes from top to bottom
    # Output: list of nodes in "right[left right[left right[...]]]" order (all right daughters are even-indexed). 

    # Important: treats specifiers as terminals, and assumes this process has already run on them. 
    # Important: assumes the strict LCA, and operates over constituency trees, not dependency trees.

# A helper to retrieve the value of a given tier on a given node. 
def val(node, tier):
    return None if node is None else node.vals[tier]

# A helper to check that a tree is well-formed. 
def validate_tree(tree):
    if len(tree) % 2 == 0: raise ValueError('tree has even number of nodes')
    if any(val(tree[-1], t) is not None for t in tier_list):
        raise ValueError('bottom right node is not empty')

# An evaluator which checks if a given tree could have been created with an MTSM or not, given a specific tau. 
def eval_tauful(tree, tau):
    validate_tree(tree)
    for i in range(0, len(tree)-1, 2):
        mother, left, right = tree[i], tree[i+1], tree[i+2]
        if any(tau(left, right).vals[t] != val(mother, t) for t in tier_list):
            return False
    return True

# An evaluator which checks if there exists a possible tau that could have created this MTSM. 
def eval_tauless(tree):
    validate_tree(tree)
    seen = {} # dict of tier, leftval, rightval -> motherval, to enforce independence
    for i in range(0, len(tree)-1, 2):
        mother, left, right = tree[i], tree[i+1], tree[i+2]
        for tier in tier_list:
            l, r, m = val(left, tier), val(right, tier), val(mother, tier)
            if m != l and m != r: 
                return False # check conservation
            if seen.setdefault((tier, l, r), m) != m: 
                return False # check independence
    return True

# An example tau, which first checks a person feature hierarchy, and then defaults to the left node's feature unless it's unvalued. (This tau is associative.)
person_hierarchy = {'1': 0, '2': 1, '3': 2}
def tau(left, right):
    mother = snode(left, right)
    for t in tier_list:
        l, r = val(left, t), val(right, t)
        if l is None or r is None:
            mother.vals[t] = r if l is None else l
        elif t == 'person' and person_hierarchy[r] < person_hierarchy[l]:
            mother.vals[t] = r
        else: 
            mother.vals[t] = l
    return mother
\end{minted}

\section{Trying it out}\label{sec:mtsm-try}

Build an input stream below and press \textbf{run}. Each row is a left daughter, from the top of the tree to the bottom; leave a tier at -- for an unvalued feature. The form starts with a third-person plural subject above a first-person singular object, so the person hierarchy lets the object's [1] win at every mother above it.

\begin{nodestream}[person=1/2/3, number=sg/pl, gender=m/f]
C
subject: person=3, number=pl
T
object: person=1, number=sg, gender=f
\end{nodestream}

\begin{minted}{python}
# stream_data comes from the form: one dict per left daughter, top to bottom
stream = []
for row in stream_data:
    n = snode()
    n.vals.update({t: row[t] for t in tier_list if row.get(t)})
    stream.append(n)

tree = run_mtsm(stream, tau)

for i, n in enumerate(tree):
    if i % 2:
        name = stream_data[i // 2]['label'] or f'left {i // 2 + 1}'
    else:
        name = 'mother' if n is not None else '(root)'
    vals = ', '.join(f'{t}={v}' for t, v in n.vals.items() if v) if n else ''
    print(f'{i:>2}  {name:<12} {vals or "–"}')

print()
print('built by this tau?   ', eval_tauful(tree, tau))
print('built by some tau?   ', eval_tauless(tree))
\end{minted}

The two evaluators answer different questions: \texttt{eval\_tauful} asks whether \emph{this} $\tau$ built the tree, while \texttt{eval\_tauless} asks whether \emph{any} conservative, tier-independent $\tau$ could have.\comment{Add a tree here once the MTSM diagrams are drawn.}
